\begin{pastexam}[label={px:ntust-115-ice-05}]{臺科大資工所 115 學年度 第 I-5 題}
（單選）

(5\%) When performing DFS or BFS traversal on a cyclic graph, what could happen if you forget to use a
visited array (or equivalent) to mark visited vertices?
\begin{enumerate}[label=(\alph*)]
\item The traversal may fall into an infinite loop, continuing until a stack overflow (for DFS) or memory
exhaustion occurs.
\item The program crashes immediately.
\item The algorithm will stop automatically, but the result will be inaccurate.
\item The program can still perform correctly.
\item There is no effect.
\item The traversal will still terminate because DFS/BFS never visits the same vertex more than once by
design.
\end{enumerate}
\end{pastexam}

\begin{pastsoln}
答案是 (a)。

DFS、BFS 不重複走同一個頂點，靠的就是 visited 標記（資結書虛擬碼~21.1、21.3 的顏色）。 拿掉之後，沿著環可以回到
走過的頂點，又從它出發再走一遍。 例：三角形 $u$--$v$--$w$--$u$，從 $u$ 做遞迴 DFS：$u \to v \to w \to u \to v \to
\cdots$，遞迴永遠不會結束，直到呼叫堆疊用完（stack overflow）。 BFS 也一樣：每取出一個頂點就把所有鄰居放進佇列，
走過的頂點一再被放回去，佇列越來越長，直到記憶體用完。

\begin{itemize}
\item[(b)] 不會「立刻」當掉，要跑到堆疊或記憶體耗盡才出事。
\item[(c)] 不會自己停下來，見上面的例子。
\item[(d)]、(e) 程式無法結束，不能算正確執行，也不是沒有影響。
\item[(f)] 「每個頂點只走一次」正是靠 visited 標記做到的，不是演算法天生的性質。
\end{itemize}
\end{pastsoln}
